\chapter{Muscle Spasm}

\section*{Definition}
A muscle spasm (muscle cramp) is a sudden, involuntary, painful contraction of a muscle or muscle group lasting seconds to minutes, often followed by residual soreness. It is distinct from dystonia, myoclonus, and tetany.

\section*{What Happens in Your Body}
Spasms result from hyperexcitability of the alpha motor neuron or the muscle fiber itself. Sustained contraction compresses local blood vessels, leading to ischemia and accumulation of metabolites, which further stimulates nociceptive fibers. Proposed mechanisms include electrolyte derangements (low sodium, potassium, magnesium, calcium), dehydration, metabolic accumulation from exercise, and reduced inhibitory input from Golgi tendon organs and Renshaw cells.

\section*{Causes}
\begin{itemize}
  \item \textbf{Physiologic:} Exercise-induced fatigue, dehydration, electrolyte loss from sweating, prolonged postures
  \item \textbf{Metabolic:} Hyponatremia, hypokalemia, hypocalcemia, hypomagnesemia, renal or liver disease
  \item \textbf{Medication-related:} Diuretics, statins, beta-agonists, loop diuretics, ACE inhibitors, niacin
  \item \textbf{Neurologic:} Motor neuron disease, radiculopathy, peripheral neuropathy, nocturnal leg cramps
  \item \textbf{Endocrine:} Hypothyroidism, hyperparathyroidism, adrenal insufficiency
  \item \textbf{Toxin-related:} Strychnine poisoning, tetanus, black widow spider venom
  \item \textbf{Vascular:} Claudication, venous insufficiency
\end{itemize}

\section*{When to See a Doctor}
See your doctor if:
\begin{itemize}
  \item Spasms are frequent, severe, or interfere with sleep or daily function
  \item Generalized cramps with muscle weakness or muscle wasting
  \item Associated with medication initiation or dose change
  \item Cramps with concurrent sensory loss, swelling, or skin changes
\end{itemize}

\section*{Related Symptoms}
\begin{itemize}
  \item Visible or palpable muscle knotting during contraction
  \item Pain localized to the contracting muscle
  \item Post-spasm tenderness lasting hours to days
  \item Fasciculations (if neuropathic cause)
  \item Nocturnal cramps in calves, feet, or hamstrings
\end{itemize}
\textbf{Important:} Nocturnal leg cramps are common in older adults and often respond to stretching the calf muscle before bedtime.

\section*{Self-Care / Home Management}
\begin{itemize}
  \item Passive stretching of the affected muscle in opposite direction of contraction
  \item Massage and application of heat (for chronic tension) or ice (for acute spasm)
  \item Hydration with electrolyte-containing fluids during and after exercise
  \item Over-the-counter magnesium supplements or tonic water (quinine not recommended due to safety concerns)
  \item Proper warm-up and cool-down routines before physical activity
\end{itemize}

\section*{How Your Doctor Will Evaluate You}
\begin{itemize}
  \item Clinical history: timing, triggers, frequency, medications, exercise habits
  \item Neuromuscular examination to assess strength, reflexes, and sensation
  \item Laboratory evaluation: electrolytes (Na, K, Ca, Mg, phosphate), renal function, TSH, CK
  \item Electromyography if neuropathic or myopathic cause is suspected
  \item Review of medications for iatrogenic causes
\end{itemize}

\section*{Treatment Options}
\begin{itemize}
  \item Treatment of underlying electrolyte imbalance or metabolic disorder
  \item Discontinuation or dose reduction of offending medications
  \item Quinine derivatives are no longer recommended due to cardiac toxicity risk
  \item Calcium channel blockers (verapamil, diltiazem) for refractory cramping in some neurologic conditions
  \item Gabapentin or pregabalin for neurogenic cramping
  \item Botulinum toxin injection for focal, refractory muscle spasms
\end{itemize}

\clearpage
